% !TeX spellcheck = en_US
\documentclass[french]{yLectureNote}

\title{Introduction à la termodynamique}
\subtitle{Thermodynamique}
\author{Paulhenry Saux}
\date{\today}
\yLanguage{Français}

\professor{}

\usepackage{graphicx}%----pour mettre des images
\usepackage[utf8]{inputenc}%---encodage
\usepackage{geometry}%---pour modifier les tailles et mettre a4paper
%\usepackage{awesomebox}%---pour les boites d'exercices, de pbq et de croquis ---d\'esactiv\'e pour les TP de PC
\usepackage{tikz}%---pour deiffner + d\'ependance de chemfig
\usepackage{tkz-tab}
\usepackage{tabularx}%---pour dimensionner automatiquement les tableaux avec variable X
\usepackage{awesomebox}%---Pour les boites info, danger et autres
\usepackage{menukeys}%---Pour deiffner les touches de Calculatrice
\usepackage{fancyhdr}%---pour les en-t\^ete personnalis\'ees
\usepackage{blindtext}%---pour les liens
\usepackage{hyperref}%---pour les liens (\`a mettre en dernier)
\usepackage{caption}%---pour la francisation de la l\'egende table vers Tableau
\usepackage{pifont}
\usepackage{array}%---pour les tableaux
\usepackage{lipsum}
\usepackage{yFlatTable}
\usepackage{multicol}
\newcommand{\Lim}[1]{\lim\limits_{\substack{#1}}\:}
\renewcommand{\vec}{\overrightarrow}
\newcommand{\bz}{\overline{z}}
\DeclareMathOperator{\card}{card}
\renewcommand{\vec}{\overrightarrow}
\newcommand{\N}[0]{\mathbb{N}}
\newcommand{\R}[0]{\mathbb{R}}
\newcommand{\C}[0]{\mathbb{C}}
\newcommand{\dd}[0]{\mathrm{d}}
\newcommand{\er}{\vec{e_{\rho}}}
\newcommand{\ep}{\vec{e_{\varphi}}}
\newcommand{\pip}{\pi^+}
\newcommand{\pim}{\pi^-}
\newcommand{\mv}{\mathcal{V}}
\DeclareMathOperator{\Div}{Div}
\newcommand{\rot}{\vec{\mathrm{rot}}}
\newcommand{\grad}{\vec{\mathrm{grad}}}
\begin{document}
\setcounter{chapter}{2}
\chapter{Principe de conservation de l'énergie}
\section{Énoncé}
\begin{theorem}[Principe de conservation de Lavoisier]
    L'énergie est une grandeur conservative, elle ne peut \^etre ni crée, ni détruite.
\end{theorem}
\begin{theorem}[Deuxième définition]
    L'énergie d'un système isolé est une constante. L'énergie de l'univers est constante.
\end{theorem}
\section{Premier principe de la thermodynamique}
\subsection{Énoncé du premier principe}

\begin{theorem}[Énoncé du premier principe]
    $$\Delta_{if} (E_c+E_p+U) = W_{i \to f} + Q_{i\to f}$$
\end{theorem}
\criticalInfo{Utilité}{On ne s'en sert que pour calculer Q}
\section{Capacités thermiques}
\begin{definition}[Capacité thermique]
    La petite quantité d'énergie, notée $c$ qu'il faut fournir sous forme de chaleur pour augmenter sa temérature de $\dd T$:
    $$\delta Q = C\dd T$$ et $$\Delta Q = c\cdot m\cdot \Delta T $$
\end{definition}
\subsection{Capacité thermique à volume constant}

\begin{theorem}[Expression de la capacité thermique volume constant]
    $C_v = (\frac{\partial U}{\partial T})_V$.
\end{theorem}
\subsection{Capacité thermique à pression constante}
\subsubsection{Introduction}
Si la transormation est monobare, on a $$P_{ext} = Cst = p_i = p_f$$
On a 
\begin{flalign*}
    Q &= \Delta U - W\\
    &=\Delta U - \int - P_{ext}\dd V\\
    &= \Delta U + P_{ext}\Delta v
    &= U_f + P_{ext}V_f - (U_i+P_{ext}V_i)\\
    &= \Delta (U+pV)
\end{flalign*}
On pose $H = U+Pv$ qui est une fonction d'état extensive, notée Enthalpie. Elle correspond à la quantité de chaleur échangée sur une transformation monobare.

Si la transformation est monobare, on a $Q = \Delta H$
\begin{theorem}[Enthalpie]
    Notée $H = U+Pv$, elle permte d'écrire $Q = \Delta H$ pour une transformation monobare.
\end{theorem}

\begin{theorem}[Expression de la capacité thermique pression constante]
    $C_p = (\frac{\partial H}{\partial T})_P$.
\end{theorem}
\section{Phénoménologie}
\subsection{Incompressible}
Le volume est constant, par exemple l'eau liquide et les solides. On a donc $V = Cst$

De plus, $\dd U = C_v \dd T + \frac{\partial U}{\partial V}\dd V \Rightarrow U = C_v T+Cst$
L'énergie interne est proportionelle à la température
\subsection{Gaz parfaits}
On a $pV = nRT$ ou $pV = mrT$ ou $ p = \rho rT$ avec $r=\frac{R}{M}$
\begin{theorem}[Première loi de Joule]
    $U = U(T)$
\end{theorem}
Donc $\dd U = C_v \dd T + \frac{\partial U}{\partial V}\dd V \Rightarrow U = C_v T + Cst$

\section{Applications aux gaz parfaits}
\subsection{Relation de Mayer}

\begin{theorem}[Relation de Mayer]
    $$C_p-C_v = nR$$
\end{theorem}
Posons maintenant $\gamma = \frac{C_p}{C_v}$. Il est connu la plupart du temps.
\begin{theorem}[Expression de Gamma]
    $$\gamma = \frac{C_P}{C_V}$$
\end{theorem}
\warningInfo{Résolution du système}{On en déduit $C_V = \frac{nR}{\gamma-1}$ et $C_p = \frac{nR\gamma}{\gamma-1}$}
\subsection{Relation de Laplace}
\criticalInfo{Hypothèses nécessaires}{
    \begin{itemize}
        \item Gaz parfait
        \item Transformation adiabatique (sans échange de chaleur avec le milieu extérieur = parois imperméables )
        \item Transformation quasi-statique
    \end{itemize}
}
\begin{theorem}[Relation de Laplace]
    $$TV^{\gamma-1} = K$$ ou $pV^\gamma = K'$ ou $p^{1-\gamma}T^\gamma=K''$.
\end{theorem}

\end{document}


